%% Figure: the cells constrained by the formulas G_4(6,10) and G_4(4,10,10) of Theorem 6.1 (ii), (iii)
%% (regions from figures.py). Each line of length n = 2,...,10 is drawn with time downwards.
\begin{tikzpicture}[x=1.85mm,y=1.85mm]
  \begin{scope}
    %% generated by figures.py; do not edit
\begin{scope}[shift={(0,0)}]
\fill[gFull] (0,0) rectangle (2,-1);
\fill[gFull] (0,-1) rectangle (2,-2);
\fill[stF] (0,-2) rectangle (2,-3);
\draw[grid,line width=0.25pt] (0,0) rectangle (2,-3);
\node[font=\scriptsize,anchor=south] at (1,0.15) {$2$};
\end{scope}
\begin{scope}[shift={(4,0)}]
\fill[gFull] (0,0) rectangle (3,-1);
\fill[gFull] (0,-1) rectangle (3,-2);
\fill[gFull] (0,-2) rectangle (3,-3);
\fill[gFull] (0,-3) rectangle (3,-4);
\fill[stF] (0,-4) rectangle (3,-5);
\draw[grid,line width=0.25pt] (0,0) rectangle (3,-5);
\node[font=\scriptsize,anchor=south] at (1.5,0.15) {$3$};
\end{scope}
\begin{scope}[shift={(9,0)}]
\fill[gFull] (0,0) rectangle (4,-1);
\fill[gFull] (0,-1) rectangle (4,-2);
\fill[gFull] (0,-2) rectangle (4,-3);
\fill[gFull] (0,-3) rectangle (4,-4);
\fill[gFull] (0,-4) rectangle (4,-5);
\fill[gFull] (0,-5) rectangle (4,-6);
\fill[stF] (0,-6) rectangle (4,-7);
\draw[grid,line width=0.25pt] (0,0) rectangle (4,-7);
\node[font=\scriptsize,anchor=south] at (2,0.15) {$4$};
\end{scope}
\begin{scope}[shift={(15,0)}]
\fill[gFull] (0,0) rectangle (5,-1);
\fill[gFull] (0,-1) rectangle (5,-2);
\fill[gFull] (0,-2) rectangle (5,-3);
\fill[gFull] (0,-3) rectangle (5,-4);
\fill[gFull] (0,-4) rectangle (5,-5);
\fill[gFull] (0,-5) rectangle (5,-6);
\fill[gFull] (0,-6) rectangle (5,-7);
\fill[gFull] (0,-7) rectangle (5,-8);
\fill[stF] (0,-8) rectangle (5,-9);
\draw[grid,line width=0.25pt] (0,0) rectangle (5,-9);
\node[font=\scriptsize,anchor=south] at (2.5,0.15) {$5$};
\end{scope}
\begin{scope}[shift={(22,0)}]
\fill[gFull] (0,0) rectangle (6,-1);
\fill[gFull] (0,-1) rectangle (6,-2);
\fill[gFull] (0,-2) rectangle (6,-3);
\fill[gFull] (0,-3) rectangle (6,-4);
\fill[gFull] (0,-4) rectangle (6,-5);
\fill[gFull] (0,-5) rectangle (6,-6);
\fill[gFull] (0,-6) rectangle (6,-7);
\fill[gFull] (0,-7) rectangle (6,-8);
\fill[gFull] (0,-8) rectangle (6,-9);
\fill[gFull] (0,-9) rectangle (6,-10);
\fill[stF] (0,-10) rectangle (6,-11);
\draw[grid,line width=0.25pt] (0,0) rectangle (6,-11);
\node[font=\scriptsize,anchor=south] at (3,0.15) {$6$};
\end{scope}
\begin{scope}[shift={(30,0)}]
\fill[gFull] (0,0) rectangle (7,-1);
\fill[gFull] (0,-1) rectangle (7,-2);
\fill[gFull] (0,-2) rectangle (7,-3);
\fill[gFull] (0,-3) rectangle (7,-4);
\fill[gFull] (0,-4) rectangle (7,-5);
\fill[gFull] (0,-5) rectangle (7,-6);
\fill[gFull] (0,-6) rectangle (7,-7);
\fill[gFull] (0,-7) rectangle (7,-8);
\fill[gFull] (0,-8) rectangle (7,-9);
\fill[gFull] (0,-9) rectangle (7,-10);
\fill[gFull] (0,-10) rectangle (7,-11);
\fill[gFull] (0,-11) rectangle (7,-12);
\fill[stF] (0,-12) rectangle (7,-13);
\draw[grid,line width=0.25pt] (0,0) rectangle (7,-13);
\node[font=\scriptsize,anchor=south] at (3.5,0.15) {$7$};
\end{scope}
\begin{scope}[shift={(39,0)}]
\fill[gEarly] (0,0) rectangle (8,-1);
\fill[gEarly] (0,-1) rectangle (8,-2);
\fill[gEarly] (0,-2) rectangle (8,-3);
\fill[gEarly] (0,-3) rectangle (8,-4);
\fill[gEarly] (0,-4) rectangle (8,-5);
\fill[gEarly] (0,-5) rectangle (8,-6);
\fill[gEarly] (0,-6) rectangle (8,-7);
\fill[gEarly] (0,-7) rectangle (7,-8);
\fill[gReturn] (7,-7) rectangle (8,-8);
\fill[gEarly] (0,-8) rectangle (6,-9);
\fill[gReturn] (6,-8) rectangle (7,-9);
\fill[gEarly] (7,-8) rectangle (8,-9);
\fill[gEarly] (0,-9) rectangle (5,-10);
\fill[gReturn] (5,-9) rectangle (6,-10);
\fill[gEarly] (6,-9) rectangle (8,-10);
\fill[gEarly] (0,-10) rectangle (4,-11);
\fill[gReturn] (4,-10) rectangle (5,-11);
\fill[gEarly] (5,-10) rectangle (8,-11);
\fill[gEarly] (0,-11) rectangle (3,-12);
\fill[gReturn] (3,-11) rectangle (4,-12);
\fill[gEarly] (4,-11) rectangle (8,-12);
\fill[gEarly] (0,-12) rectangle (2,-13);
\fill[gReturn] (2,-12) rectangle (3,-13);
\fill[gEarly] (3,-12) rectangle (8,-13);
\fill[gEarly] (0,-13) rectangle (1,-14);
\fill[gReturn] (1,-13) rectangle (2,-14);
\fill[gEarly] (2,-13) rectangle (8,-14);
\fill[gNone] (0,-14) rectangle (8,-15);
\draw[grid,line width=0.25pt] (0,0) rectangle (8,-15);
\node[font=\scriptsize,anchor=south] at (4,0.15) {$8$};
\end{scope}
\begin{scope}[shift={(49,0)}]
\fill[gEarly] (0,0) rectangle (9,-1);
\fill[gEarly] (0,-1) rectangle (9,-2);
\fill[gEarly] (0,-2) rectangle (9,-3);
\fill[gEarly] (0,-3) rectangle (9,-4);
\fill[gEarly] (0,-4) rectangle (9,-5);
\fill[gEarly] (0,-5) rectangle (9,-6);
\fill[gEarly] (0,-6) rectangle (9,-7);
\fill[gEarly] (0,-7) rectangle (9,-8);
\fill[gEarly] (0,-8) rectangle (8,-9);
\fill[gReturn] (8,-8) rectangle (9,-9);
\fill[gEarly] (0,-9) rectangle (7,-10);
\fill[gReturn] (7,-9) rectangle (8,-10);
\fill[gEarly] (8,-9) rectangle (9,-10);
\fill[gEarly] (0,-10) rectangle (6,-11);
\fill[gReturn] (6,-10) rectangle (7,-11);
\fill[gEarly] (7,-10) rectangle (9,-11);
\fill[gEarly] (0,-11) rectangle (5,-12);
\fill[gReturn] (5,-11) rectangle (6,-12);
\fill[gEarly] (6,-11) rectangle (9,-12);
\fill[gEarly] (0,-12) rectangle (4,-13);
\fill[gReturn] (4,-12) rectangle (5,-13);
\fill[gEarly] (5,-12) rectangle (9,-13);
\fill[gEarly] (0,-13) rectangle (3,-14);
\fill[gReturn] (3,-13) rectangle (4,-14);
\fill[gEarly] (4,-13) rectangle (9,-14);
\fill[gEarly] (0,-14) rectangle (2,-15);
\fill[gReturn] (2,-14) rectangle (3,-15);
\fill[gEarly] (3,-14) rectangle (9,-15);
\fill[gNone] (0,-15) rectangle (9,-16);
\fill[gNone] (0,-16) rectangle (9,-17);
\draw[grid,line width=0.25pt] (0,0) rectangle (9,-17);
\node[font=\scriptsize,anchor=south] at (4.5,0.15) {$9$};
\end{scope}
\begin{scope}[shift={(60,0)}]
\fill[gEarly] (0,0) rectangle (10,-1);
\fill[gEarly] (0,-1) rectangle (10,-2);
\fill[gEarly] (0,-2) rectangle (10,-3);
\fill[gEarly] (0,-3) rectangle (10,-4);
\fill[gEarly] (0,-4) rectangle (10,-5);
\fill[gEarly] (0,-5) rectangle (10,-6);
\fill[gEarly] (0,-6) rectangle (10,-7);
\fill[gEarly] (0,-7) rectangle (10,-8);
\fill[gEarly] (0,-8) rectangle (10,-9);
\fill[gEarly] (0,-9) rectangle (9,-10);
\fill[gReturn] (9,-9) rectangle (10,-10);
\fill[gEarly] (0,-10) rectangle (8,-11);
\fill[gReturn] (8,-10) rectangle (9,-11);
\fill[gEarly] (9,-10) rectangle (10,-11);
\fill[gEarly] (0,-11) rectangle (7,-12);
\fill[gReturn] (7,-11) rectangle (8,-12);
\fill[gEarly] (8,-11) rectangle (10,-12);
\fill[gEarly] (0,-12) rectangle (6,-13);
\fill[gReturn] (6,-12) rectangle (7,-13);
\fill[gEarly] (7,-12) rectangle (10,-13);
\fill[gEarly] (0,-13) rectangle (5,-14);
\fill[gReturn] (5,-13) rectangle (6,-14);
\fill[gEarly] (6,-13) rectangle (10,-14);
\fill[gEarly] (0,-14) rectangle (4,-15);
\fill[gReturn] (4,-14) rectangle (5,-15);
\fill[gEarly] (5,-14) rectangle (10,-15);
\fill[gEarly] (0,-15) rectangle (3,-16);
\fill[gReturn] (3,-15) rectangle (4,-16);
\fill[gEarly] (4,-15) rectangle (10,-16);
\fill[gNone] (0,-16) rectangle (10,-17);
\fill[gNone] (0,-17) rectangle (10,-18);
\fill[gNone] (0,-18) rectangle (10,-19);
\draw[grid,line width=0.25pt] (0,0) rectangle (10,-19);
\node[font=\scriptsize,anchor=south] at (5,0.15) {$10$};
\end{scope}

    \node[font=\small,anchor=south west] at (-0.5,2.6) {(a) $\G_4(6,10)$, \cref{thm:four}~(ii): $T=6$, $N_1=10$};
  \end{scope}
  \begin{scope}[shift={(0,-25.5)}]
    %% generated by figures.py; do not edit
\begin{scope}[shift={(0,0)}]
\fill[gFull] (0,0) rectangle (2,-1);
\fill[gFull] (0,-1) rectangle (2,-2);
\fill[stF] (0,-2) rectangle (2,-3);
\draw[grid,line width=0.25pt] (0,0) rectangle (2,-3);
\node[font=\scriptsize,anchor=south] at (1,0.15) {$2$};
\end{scope}
\begin{scope}[shift={(4,0)}]
\fill[gFull] (0,0) rectangle (3,-1);
\fill[gFull] (0,-1) rectangle (3,-2);
\fill[gFull] (0,-2) rectangle (3,-3);
\fill[gFull] (0,-3) rectangle (3,-4);
\fill[stF] (0,-4) rectangle (3,-5);
\draw[grid,line width=0.25pt] (0,0) rectangle (3,-5);
\node[font=\scriptsize,anchor=south] at (1.5,0.15) {$3$};
\end{scope}
\begin{scope}[shift={(9,0)}]
\fill[gFull] (0,0) rectangle (4,-1);
\fill[gFull] (0,-1) rectangle (4,-2);
\fill[gFull] (0,-2) rectangle (4,-3);
\fill[gFull] (0,-3) rectangle (4,-4);
\fill[gFull] (0,-4) rectangle (4,-5);
\fill[gFull] (0,-5) rectangle (4,-6);
\fill[stF] (0,-6) rectangle (4,-7);
\draw[grid,line width=0.25pt] (0,0) rectangle (4,-7);
\node[font=\scriptsize,anchor=south] at (2,0.15) {$4$};
\end{scope}
\begin{scope}[shift={(15,0)}]
\fill[gFull] (0,0) rectangle (5,-1);
\fill[gFull] (0,-1) rectangle (5,-2);
\fill[gFull] (0,-2) rectangle (5,-3);
\fill[gFull] (0,-3) rectangle (5,-4);
\fill[gFull] (0,-4) rectangle (5,-5);
\fill[gFull] (0,-5) rectangle (5,-6);
\fill[gFull] (0,-6) rectangle (5,-7);
\fill[gFull] (0,-7) rectangle (5,-8);
\fill[stF] (0,-8) rectangle (5,-9);
\draw[grid,line width=0.25pt] (0,0) rectangle (5,-9);
\node[font=\scriptsize,anchor=south] at (2.5,0.15) {$5$};
\end{scope}
\begin{scope}[shift={(22,0)}]
\fill[gEarly] (0,0) rectangle (6,-1);
\fill[gEarly] (0,-1) rectangle (6,-2);
\fill[gEarly] (0,-2) rectangle (6,-3);
\fill[gEarly] (0,-3) rectangle (6,-4);
\fill[gEarly] (0,-4) rectangle (6,-5);
\fill[gEarly] (0,-5) rectangle (5,-6);
\fill[gReturn] (5,-5) rectangle (6,-6);
\fill[gEarly] (0,-6) rectangle (4,-7);
\fill[gReturn] (4,-6) rectangle (5,-7);
\fill[gEarly] (5,-6) rectangle (6,-7);
\fill[gEarly] (0,-7) rectangle (3,-8);
\fill[gReturn] (3,-7) rectangle (4,-8);
\fill[gEarly] (4,-7) rectangle (6,-8);
\fill[gEarly] (0,-8) rectangle (2,-9);
\fill[gReturn] (2,-8) rectangle (3,-9);
\fill[gEarly] (3,-8) rectangle (6,-9);
\fill[gEarly] (0,-9) rectangle (1,-10);
\fill[gReturn] (1,-9) rectangle (2,-10);
\fill[gEarly] (2,-9) rectangle (6,-10);
\fill[gNone] (0,-10) rectangle (5,-11);
\fill[stF] (5,-10) rectangle (6,-11);
\draw[grid,line width=0.25pt] (0,0) rectangle (6,-11);
\node[font=\scriptsize,anchor=south] at (3,0.15) {$6$};
\end{scope}
\begin{scope}[shift={(30,0)}]
\fill[gEarly] (0,0) rectangle (7,-1);
\fill[gEarly] (0,-1) rectangle (7,-2);
\fill[gEarly] (0,-2) rectangle (7,-3);
\fill[gEarly] (0,-3) rectangle (7,-4);
\fill[gEarly] (0,-4) rectangle (7,-5);
\fill[gEarly] (0,-5) rectangle (7,-6);
\fill[gEarly] (0,-6) rectangle (6,-7);
\fill[gReturn] (6,-6) rectangle (7,-7);
\fill[gEarly] (0,-7) rectangle (5,-8);
\fill[gReturn] (5,-7) rectangle (6,-8);
\fill[gEarly] (6,-7) rectangle (7,-8);
\fill[gEarly] (0,-8) rectangle (4,-9);
\fill[gReturn] (4,-8) rectangle (5,-9);
\fill[gEarly] (5,-8) rectangle (7,-9);
\fill[gEarly] (0,-9) rectangle (3,-10);
\fill[gReturn] (3,-9) rectangle (4,-10);
\fill[gEarly] (4,-9) rectangle (7,-10);
\fill[gEarly] (0,-10) rectangle (2,-11);
\fill[gReturn] (2,-10) rectangle (3,-11);
\fill[gEarly] (3,-10) rectangle (7,-11);
\fill[gNone] (0,-11) rectangle (5,-12);
\fill[gCone] (5,-11) rectangle (7,-12);
\fill[gNone] (0,-12) rectangle (6,-13);
\fill[stF] (6,-12) rectangle (7,-13);
\draw[grid,line width=0.25pt] (0,0) rectangle (7,-13);
\node[font=\scriptsize,anchor=south] at (3.5,0.15) {$7$};
\end{scope}
\begin{scope}[shift={(39,0)}]
\fill[gEarly] (0,0) rectangle (8,-1);
\fill[gEarly] (0,-1) rectangle (8,-2);
\fill[gEarly] (0,-2) rectangle (8,-3);
\fill[gEarly] (0,-3) rectangle (8,-4);
\fill[gEarly] (0,-4) rectangle (8,-5);
\fill[gEarly] (0,-5) rectangle (8,-6);
\fill[gEarly] (0,-6) rectangle (8,-7);
\fill[gEarly] (0,-7) rectangle (7,-8);
\fill[gReturn] (7,-7) rectangle (8,-8);
\fill[gEarly] (0,-8) rectangle (6,-9);
\fill[gReturn] (6,-8) rectangle (7,-9);
\fill[gEarly] (7,-8) rectangle (8,-9);
\fill[gEarly] (0,-9) rectangle (5,-10);
\fill[gReturn] (5,-9) rectangle (6,-10);
\fill[gEarly] (6,-9) rectangle (8,-10);
\fill[gEarly] (0,-10) rectangle (4,-11);
\fill[gReturn] (4,-10) rectangle (5,-11);
\fill[gEarly] (5,-10) rectangle (8,-11);
\fill[gEarly] (0,-11) rectangle (3,-12);
\fill[gReturn] (3,-11) rectangle (4,-12);
\fill[gEarly] (4,-11) rectangle (8,-12);
\fill[gNone] (0,-12) rectangle (5,-13);
\fill[gCone] (5,-12) rectangle (8,-13);
\fill[gNone] (0,-13) rectangle (6,-14);
\fill[gCone] (6,-13) rectangle (8,-14);
\fill[gNone] (0,-14) rectangle (7,-15);
\fill[stF] (7,-14) rectangle (8,-15);
\draw[grid,line width=0.25pt] (0,0) rectangle (8,-15);
\node[font=\scriptsize,anchor=south] at (4,0.15) {$8$};
\end{scope}
\begin{scope}[shift={(49,0)}]
\fill[gEarly] (0,0) rectangle (9,-1);
\fill[gEarly] (0,-1) rectangle (9,-2);
\fill[gEarly] (0,-2) rectangle (9,-3);
\fill[gEarly] (0,-3) rectangle (9,-4);
\fill[gEarly] (0,-4) rectangle (9,-5);
\fill[gEarly] (0,-5) rectangle (9,-6);
\fill[gEarly] (0,-6) rectangle (9,-7);
\fill[gEarly] (0,-7) rectangle (9,-8);
\fill[gEarly] (0,-8) rectangle (8,-9);
\fill[gReturn] (8,-8) rectangle (9,-9);
\fill[gEarly] (0,-9) rectangle (7,-10);
\fill[gReturn] (7,-9) rectangle (8,-10);
\fill[gEarly] (8,-9) rectangle (9,-10);
\fill[gEarly] (0,-10) rectangle (6,-11);
\fill[gReturn] (6,-10) rectangle (7,-11);
\fill[gEarly] (7,-10) rectangle (9,-11);
\fill[gEarly] (0,-11) rectangle (5,-12);
\fill[gReturn] (5,-11) rectangle (6,-12);
\fill[gEarly] (6,-11) rectangle (9,-12);
\fill[gEarly] (0,-12) rectangle (4,-13);
\fill[gReturn] (4,-12) rectangle (5,-13);
\fill[gEarly] (5,-12) rectangle (9,-13);
\fill[gNone] (0,-13) rectangle (5,-14);
\fill[gCone] (5,-13) rectangle (9,-14);
\fill[gNone] (0,-14) rectangle (6,-15);
\fill[gCone] (6,-14) rectangle (9,-15);
\fill[gNone] (0,-15) rectangle (7,-16);
\fill[gCone] (7,-15) rectangle (9,-16);
\fill[gNone] (0,-16) rectangle (8,-17);
\fill[stF] (8,-16) rectangle (9,-17);
\draw[grid,line width=0.25pt] (0,0) rectangle (9,-17);
\node[font=\scriptsize,anchor=south] at (4.5,0.15) {$9$};
\end{scope}
\begin{scope}[shift={(60,0)}]
\fill[gEarly] (0,0) rectangle (10,-1);
\fill[gEarly] (0,-1) rectangle (10,-2);
\fill[gEarly] (0,-2) rectangle (10,-3);
\fill[gEarly] (0,-3) rectangle (10,-4);
\fill[gEarly] (0,-4) rectangle (10,-5);
\fill[gEarly] (0,-5) rectangle (10,-6);
\fill[gEarly] (0,-6) rectangle (10,-7);
\fill[gEarly] (0,-7) rectangle (10,-8);
\fill[gEarly] (0,-8) rectangle (10,-9);
\fill[gEarly] (0,-9) rectangle (9,-10);
\fill[gReturn] (9,-9) rectangle (10,-10);
\fill[gEarly] (0,-10) rectangle (8,-11);
\fill[gReturn] (8,-10) rectangle (9,-11);
\fill[gEarly] (9,-10) rectangle (10,-11);
\fill[gEarly] (0,-11) rectangle (7,-12);
\fill[gReturn] (7,-11) rectangle (8,-12);
\fill[gEarly] (8,-11) rectangle (10,-12);
\fill[gEarly] (0,-12) rectangle (6,-13);
\fill[gReturn] (6,-12) rectangle (7,-13);
\fill[gEarly] (7,-12) rectangle (10,-13);
\fill[gEarly] (0,-13) rectangle (5,-14);
\fill[gReturn] (5,-13) rectangle (6,-14);
\fill[gEarly] (6,-13) rectangle (10,-14);
\fill[gNone] (0,-14) rectangle (5,-15);
\fill[gCone] (5,-14) rectangle (10,-15);
\fill[gNone] (0,-15) rectangle (6,-16);
\fill[gCone] (6,-15) rectangle (10,-16);
\fill[gNone] (0,-16) rectangle (7,-17);
\fill[gCone] (7,-16) rectangle (10,-17);
\fill[gNone] (0,-17) rectangle (8,-18);
\fill[gCone] (8,-17) rectangle (10,-18);
\fill[gNone] (0,-18) rectangle (9,-19);
\fill[stF] (9,-18) rectangle (10,-19);
\draw[grid,line width=0.25pt] (0,0) rectangle (10,-19);
\node[font=\scriptsize,anchor=south] at (5,0.15) {$10$};
\end{scope}

    \node[font=\small,anchor=south west] at (-0.5,2.6) {(b) $\G_4(4,10,10)$, \cref{thm:four}~(iii): $T=4$, $N_1=C=10$};
  \end{scope}
  \node[font=\scriptsize,anchor=east] at (-0.4,-0.5) {$t=0$};
  \node[font=\scriptsize,anchor=east] at (-0.4,-26.0) {$t=0$};
  \node[anchor=north west,font=\scriptsize,align=left] at (0,-46.6) {%
    \swatch{gFull}\,(G1): the line is synchronized at time $2n-2$\qquad
    \swatch{stF}\,$\Fq$ at time $2n-2$ (G1), (G5)\\[0.4ex]
    \swatch{gEarly}\,(G2): no cell fires up to the relative time $T$\qquad
    \swatch{gReturn}\,(G3), (G5): the return chain differs from the half-line\\[0.4ex]
    \swatch{gCone}\,(G5): no cell of the cone of the right end fires\qquad
    \swatch{gNone}\,not constrained};
\end{tikzpicture}
